\chapter*{Préface}
\addcontentsline{toc}{chapter}{Préface}

Ce manuel accompagne une année de mathématiques en classe de Seconde (T10).
Il peut être utilisé intégralement ou comme une banque de leçons, d'activités,
d'exercices et d'évaluations. La progression proposée reste adaptable au rythme
de la classe et au temps réellement disponible.

L'ouvrage vise deux usages complémentaires : fournir à chaque élève un cours
lisible et progressif, et donner à chaque enseignant un ensemble directement
exploitable de situations de recherche, d'exemples résolus, d'exercices gradués,
de problèmes, d'automatismes et d'évaluations. Les prolongements destinés aux
élèves ambitieux sont clairement signalés comme extensions facultatives.

\begin{extensiont10box}[À lire avant d'utiliser le manuel]
Certaines pages proposent volontairement des notions qui ne font pas partie du
programme T10. Elles sont destinées aux élèves qui avancent beaucoup plus vite
que leur groupe, ou qui préparent un objectif sensiblement plus ambitieux
(concours, olympiades, études scientifiques exigeantes). Elles ne sont \textbf{ni
à enseigner comme cours obligatoire, ni prérequis, ni évaluables} en T10.

Le repère est volontairement sans ambiguïté : chaque défi et chaque bloc
d'approfondissement hors programme porte une \textbf{ligne rouge continue dans
la marge gauche}, un fond rouge très pâle et la mention
\textbf{« Hors programme T10 -- non exigible »}. Lorsqu'une activité ou un
exercice isolé dépasse le programme, son titre comporte aussi le repère rouge
\ExtensionTTen.\end{extensiont10box}

Les contextes ont été choisis pour ne pas supposer un accès permanent à
Internet ou à un ordinateur. Les activités numériques proposent donc toujours
une exploitation possible sur papier; Python ou GeoGebra restent des compléments.
